% Mathématiques 6e — cours V3
% Probabilités

\section{Objectifs}

\begin{itemize}
\item Comprendre qu'une expérience aléatoire peut avoir plusieurs résultats possibles.
\item Comprendre qu'une probabilité est un nombre compris entre $0$ et $1$.
\item Relier des expressions comme « une chance sur deux » à une fraction, une écriture décimale ou un pourcentage.
\item Calculer une probabilité dans des situations simples d'équiprobabilité.
\item Répéter une expérience aléatoire et calculer une fréquence observée.
\item Comparer une fréquence expérimentale à une probabilité calculée.
\item Comprendre qu'une courte série d'expériences peut s'écarter fortement du modèle théorique.
\end{itemize}

\section{1. Une expérience dont le résultat n'est pas connu à l'avance}

Lorsqu'on lance une pièce, on ne sait pas avec certitude quel côté apparaîtra.

Lorsqu'on lance un dé, on ne connaît pas à l'avance le nombre obtenu.

On parle d'\textbf{expérience aléatoire} lorsque le résultat précis ne peut pas être prévu avec certitude avant de réaliser l'expérience.

\begin{remarque}
Le vocabulaire des probabilités peut être utilisé par le professeur, mais en 6e l'essentiel est surtout de comprendre les situations, les nombres associés et le lien avec des répétitions expérimentales.
\end{remarque}

\section{2. Les résultats possibles}

Pour un dé équilibré à six faces numérotées, les résultats possibles sont :
\[
1,\quad2,\quad3,\quad4,\quad5,\quad6.
\]

Pour une pièce, on peut obtenir deux résultats :
\[
\text{pile}\qquad\text{ou}\qquad\text{face}.
\]

Avant de calculer une probabilité, il faut identifier correctement tous les résultats possibles.

\begin{attention}
Oublier un résultat possible ou compter deux fois le même résultat fausse tout le raisonnement.
\end{attention}

\section{3. Un nombre entre zéro et un}

Une probabilité est un nombre compris entre :
\[
0
\]
et
\[
1.
\]

On écrit :
\[
\boxed{0\le P\le1}.
\]

Une probabilité égale à $0$ correspond à une situation impossible.

Une probabilité égale à $1$ correspond à une situation certaine.

Entre les deux, l'évènement est possible mais non certain.

\coursefigure{/images/mathematiques/6eme/probabilites-echelle.svg}{Échelle de probabilité allant de l'impossible au certain.}{Une probabilité est toujours comprise entre $0$ et $1$ ; la position sur l'échelle traduit le degré de possibilité.}

\section{4. Différentes écritures d'une même probabilité}

Une probabilité peut être écrite sous plusieurs formes.

Par exemple :
\[
\dfrac12=0{,}5=50\%.
\]

Ces trois écritures représentent le même nombre.

De même :
\[
\dfrac14=0{,}25=25\%.
\]

Et :
\[
\dfrac34=0{,}75=75\%.
\]

\begin{propriete}
Une probabilité peut être représentée par une fraction, une écriture décimale ou un pourcentage, à condition que toutes ces écritures représentent un nombre compris entre $0$ et $1$.
\end{propriete}

\section{5. « Une chance sur deux »}

L'expression :
\[
\text{« une chance sur deux »}
\]
correspond naturellement au nombre :
\[
\dfrac12.
\]

On peut aussi écrire :
\[
0{,}5
\]
ou
\[
50\%.
\]

Cette traduction entre langage courant et nombres est un objectif important en 6e.

\section{6. Équiprobabilité}

Dans certaines expériences, tous les résultats possibles ont la même probabilité.

On parle d'\textbf{équiprobabilité}.

Un dé équilibré est modélisé en considérant ses six faces comme également probables.

Une pièce équilibrée est modélisée en considérant pile et face comme également probables.

Dans ces situations simples :
\[
P=\dfrac{\text{nombre de résultats favorables}}{\text{nombre total de résultats possibles}}.
\]

Cette formule ne doit être utilisée que lorsque les résultats sont bien équiprobables.

\section{7. Exemple développé : obtenir un nombre pair}

On lance un dé équilibré à six faces.

Les résultats possibles sont :
\[
1,\quad2,\quad3,\quad4,\quad5,\quad6.
\]

Les résultats favorables à « obtenir un nombre pair » sont :
\[
2,\quad4,\quad6.
\]

Il y a donc :
\[
3
\]
résultats favorables parmi :
\[
6
\]
résultats équiprobables.

La probabilité vaut :
\[
P(\text{pair})=\dfrac36=\dfrac12.
\]

Donc :
\[
\boxed{P(\text{pair})=0{,}5=50\%}.
\]

\section{8. Exemple développé : obtenir un nombre supérieur à quatre}

Avec le même dé, les résultats favorables sont :
\[
5,\quad6.
\]

Il y a donc :
\[
2
\]
résultats favorables sur :
\[
6.
\]

Ainsi :
\[
P(\text{supérieur à }4)=\dfrac26=\dfrac13.
\]

L'écriture fractionnaire :
\[
\dfrac13
\]
est ici exacte.

\begin{remarque}
Il n'est pas toujours nécessaire de convertir une probabilité en écriture décimale. Une fraction peut être l'écriture exacte la plus simple.
\end{remarque}

\section{9. Évènement impossible}

Avec un dé numéroté de $1$ à $6$, obtenir :
\[
8
\]
est impossible.

Aucun résultat possible n'est favorable.

On peut écrire :
\[
P(\text{obtenir }8)=\dfrac06=0.
\]

\section{10. Évènement certain}

Avec le même dé, obtenir un nombre inférieur à :
\[
7
\]
est certain.

Les six résultats possibles conviennent.

On écrit :
\[
P(\text{nombre inférieur à }7)=\dfrac66=1.
\]

\section{11. Comparer des probabilités}

Comparer deux probabilités revient à comparer deux nombres.

Par exemple :
\[
\dfrac14<\dfrac12.
\]

Donc un résultat de probabilité :
\[
\dfrac12
\]
est plus probable qu'un résultat de probabilité :
\[
\dfrac14.
\]

Cela ne signifie pas qu'il se produira obligatoirement lors de la prochaine expérience.

\begin{attention}
« Plus probable » ne signifie pas « certain ».
\end{attention}

\section{12. Répéter une expérience}

Une expérience aléatoire peut être répétée plusieurs fois.

On peut par exemple lancer une pièce :
\[
10,\quad50,\quad100,\quad500
\]
fois.

À chaque série, on compte combien de fois apparaît un résultat choisi.

Ce comptage conduit à une \textbf{fréquence observée}.

\section{13. Fréquence expérimentale}

Si un résultat apparaît $n_A$ fois au cours de $n$ répétitions, sa fréquence est :
\[
f=\dfrac{n_A}{n}.
\]

\begin{exemple}
Une pièce est lancée $20$ fois et on observe pile $12$ fois.

La fréquence de pile vaut :
\[
f=\dfrac{12}{20}=\dfrac35=0{,}6.
\]

On peut aussi écrire :
\[
f=60\%.
\]
\end{exemple}

Cette fréquence est une observation réelle sur une série d'expériences.

\section{14. Probabilité et fréquence ne sont pas la même chose}

Pour une pièce équilibrée, le modèle donne :
\[
P(\text{pile})=\dfrac12=0{,}5.
\]

Une série de $20$ lancers peut pourtant donner :
\[
f(\text{pile})=0{,}6.
\]

Il n'y a pas de contradiction.

La probabilité appartient au \textbf{modèle théorique}.

La fréquence appartient à la \textbf{série effectivement observée}.

\section{15. Les fluctuations}

Lors d'une petite série, la fréquence peut varier fortement.

Sur quatre lancers d'une pièce, on peut observer :
\[
4
\]
fois pile.

La fréquence vaut alors :
\[
1.
\]

Pourtant, le modèle d'une pièce équilibrée reste :
\[
P(\text{pile})=\dfrac12.
\]

Une courte série n'a aucune obligation de reproduire exactement la probabilité théorique.

\section{16. Que se passe-t-il lorsque l'on répète davantage ?}

Lorsque le nombre de répétitions augmente, les fréquences observées tendent souvent à devenir moins instables autour de la valeur prévue par le modèle.

\coursefigure{/images/mathematiques/6eme/probabilites-frequence.svg}{Graphique schématique de fréquences observées qui fluctuent autour d'une référence théorique.}{Les premières fréquences peuvent beaucoup varier ; avec davantage de répétitions, les fluctuations deviennent généralement moins fortes autour de la valeur théorique.}

\begin{attention}
Cela ne signifie pas que la fréquence devient obligatoirement exactement égale à la probabilité après un nombre fixé de répétitions.
\end{attention}

\section{17. Exemple développé : cinquante lancers}

On lance une pièce équilibrée :
\[
50
\]
fois.

On observe pile :
\[
27
\]
fois.

La fréquence vaut :
\[
f=\dfrac{27}{50}=0{,}54.
\]

La probabilité théorique vaut :
\[
p=0{,}5.
\]

L'écart numérique est :
\[
0{,}54-0{,}5=0{,}04.
\]

La fréquence observée est donc proche du modèle, sans lui être exactement égale.

\section{18. Exemple développé : deux séries différentes}

Deux groupes lancent chacun la même pièce équilibrée $20$ fois.

Le groupe A obtient pile :
\[
8
\]
fois.

Sa fréquence est :
\[
f_A=\dfrac8{20}=0{,}4.
\]

Le groupe B obtient pile :
\[
13
\]
fois.

Sa fréquence est :
\[
f_B=\dfrac{13}{20}=0{,}65.
\]

Les deux groupes ont utilisé le même modèle théorique :
\[
p=0{,}5.
\]

Des résultats expérimentaux différents sont donc parfaitement possibles.

\section{19. Méthode : calculer en équiprobabilité}

\begin{methode}
Dans une situation simple d'équiprobabilité :
\begin{enumerate}
\item lister tous les résultats possibles ;
\item vérifier qu'ils sont supposés également probables ;
\item repérer ceux qui réalisent la condition demandée ;
\item compter les résultats favorables ;
\item former la fraction
\[
\dfrac{\text{favorables}}{\text{possibles}};
\]
\item simplifier si possible ;
\item vérifier que le résultat appartient à l'intervalle $[0;1]$.
\end{enumerate}
\end{methode}

\section{20. Méthode : exploiter une expérience répétée}

\begin{methode}
Après plusieurs répétitions :
\begin{enumerate}
\item compter le nombre total d'essais ;
\item compter le nombre de réalisations du résultat étudié ;
\item calculer la fréquence ;
\item comparer cette fréquence à la probabilité théorique si elle est connue ;
\item décrire l'écart sans déclarer le modèle faux au seul motif que les deux nombres ne sont pas identiques.
\end{enumerate}
\end{methode}

\section{21. Un modèle doit correspondre à la situation}

La formule :
\[
\dfrac{\text{favorables}}{\text{possibles}}
\]
avec simple comptage suppose une équiprobabilité.

Elle n'est pas automatiquement valable pour n'importe quel objet.

\begin{exemple}
Une roue possède deux secteurs, mais l'un occupe les trois quarts du disque et l'autre seulement un quart.

Il y a bien deux couleurs possibles, mais elles ne sont pas équiprobables.

On ne peut donc pas conclure :
\[
\dfrac12
\]
pour chaque couleur uniquement parce qu'il existe deux couleurs.
\end{exemple}

\section{22. Contrôler un résultat}

Une probabilité de :
\[
1{,}4
\]
est impossible car :
\[
1{,}4>1.
\]

Une probabilité de :
\[
-0{,}2
\]
est également impossible car :
\[
-0{,}2<0.
\]

Le contrôle :
\[
0\le P\le1
\]
permet donc de détecter certaines erreurs immédiatement.

\section{23. Erreurs fréquentes}

\begin{itemize}
\item Croire qu'une probabilité de $\dfrac12$ impose une alternance parfaite des résultats.
\item Confondre probabilité théorique et fréquence observée.
\item Utiliser l'équiprobabilité sans vérifier que les résultats ont réellement la même chance.
\item Oublier un résultat possible.
\item Compter deux fois un même résultat favorable.
\item Produire une probabilité supérieure à $1$ sans contrôler le résultat.
\item Déclarer une pièce « truquée » après seulement quelques lancers différents de $50\%$.
\end{itemize}

\section{À retenir}

\begin{aretenir}
Une probabilité est un nombre compris entre $0$ et $1$ :
\[
\boxed{0\le P\le1}.
\]

Dans une situation simple où tous les résultats sont équiprobables :
\[
\boxed{
P=\dfrac{\text{nombre de résultats favorables}}{\text{nombre total de résultats possibles}}
}.
\]

Une fréquence expérimentale est calculée à partir de résultats réellement observés :
\[
\boxed{
f=\dfrac{\text{nombre de réalisations observées}}{\text{nombre total d'essais}}
}.
\]

La fréquence peut différer de la probabilité théorique, surtout lorsque le nombre de répétitions est faible.

Répéter davantage l'expérience permet de confronter les observations au modèle sans attendre une égalité parfaite à chaque série.
\end{aretenir}
